\documentclass{article}
\usepackage{tcolorbox}
\usepackage{amsmath}
\begin{document}

\begin{tcolorbox}[title=Box 1: Mathematical model of diverse misinformation]
We wish to explore the potential impacts of our results on the spread of diverse misinformation on social networks. We consider that multiple independent streams of misinformation spread simultaneously.

We account for the above using three stylized patterns for the network structure. First, we divide the network into two demographic classes of equal size, simply labeled 1 and 2.

According to the above, we can write the fraction of nodes $p^1_{k,\ell}$ which are of demographic class 1 with $k$ contacts of class 1 and $\ell$ contacts of class 2:
\begin{equation}
    p^1_{k,\ell} \propto \frac{1}{2}(k+\ell)^{-\alpha}\left[\frac{1}{2}\binom{k+\ell}{k}Q^k(1-Q)^\ell + \frac{1}{2}\binom{k+\ell}{k}(1/2)^{k+\ell}\right]
\end{equation}

The fraction of individuals of a certain type $(i,k,\ell)$ that are duped, $D^i_{k,\ell}$, can be followed in time:
\begin{equation}
    \frac{d}{dt} D^i_{k,\ell} = \lambda_i\left(p^i_{k,\ell}-D^i_{k,\ell}\right)\left(k\theta_{i,1}+\ell\theta_{i,2}\right)-\gamma D^i_{k,\ell}\left(k\phi_{i,1}+\ell\phi_{i,2}\right)
\end{equation}
\end{tcolorbox}

\end{document}
